\documentclass[11pt,aspectratio=169]{beamer}

\usepackage[T1]{fontenc}
\usepackage{lmodern}
\usepackage{graphicx}
\usepackage{amsmath,amssymb,mathtools}
\usepackage{xcolor}

\setbeamertemplate{navigation symbols}{}
\setbeamertemplate{itemize items}[circle]
\setbeamertemplate{footline}[frame number]
\setbeamersize{text margin left=0.65cm,text margin right=0.65cm}
\linespread{1.12}

\definecolor{mainblue}{RGB}{38,82,122}
\definecolor{mainred}{RGB}{176,73,61}

\title{Alternative Views of the Toy-Model Transition}
\subtitle{Figures for discussion}
\author{Tommaso De Santo}
\date{August 2026}

\begin{document}

{
\setbeamertemplate{footline}{}
\begin{frame}[plain,c]
  \titlepage
\end{frame}
}

% ---------------------------------------------------------------------------
\begin{frame}{Two ways to organize the same feedback}
\centering
\includegraphics[width=0.96\textwidth,height=0.67\textheight,keepaspectratio]{../output/model/fertility_population_housing_transition_note/fertility_feedback_two_interpretations.pdf}

\vspace{-0.3em}
\small
\textbf{Impact timing:} $I$ is the market-clearing impact point with inherited cohorts.

\textbf{Feedback decomposition:} $D$ holds $p=p_0$; the move to $A'$ shows how falling housing costs prevent the no-feedback demographic contraction from continuing.
\end{frame}

% ---------------------------------------------------------------------------
\begin{frame}{Alternative 1: the primitive fertility shock}
\begin{columns}[T,onlytextwidth]
\begin{column}{0.64\textwidth}
\centering
\includegraphics[width=\textwidth,height=0.68\textheight,keepaspectratio]{../output/model/fertility_population_housing_transition_note/primitive_fertility_shock.pdf}
\end{column}
\begin{column}{0.32\textwidth}
\small
\begin{itemize}
  \setlength\itemsep{1.0em}
  \item \textbf{Primitive:} $v_0\rightarrow v_1<v_0$.
  \item \textbf{Hold $p=p_0$:} fertility falls from $1$ to $n^D$.
  \item \textbf{Point $D$:} prices and cohort sizes have not yet adjusted.
\end{itemize}
\end{column}
\end{columns}
\end{frame}

% ---------------------------------------------------------------------------
\begin{frame}{Alternative 2: fast equilibrium feedback}
\centering
\includegraphics[width=0.96\textwidth,height=0.67\textheight,keepaspectratio]{../output/model/fertility_population_housing_transition_note/contemporaneous_housing_feedback.pdf}

\vspace{-0.3em}
\small
Holding $n=n^D$ gives the mechanical price $p^{PE}$. At that price households want $n^F>n^D$, so child-related demand pushes the price back to $p^I$ and fertility down to $n^I$:
\[
p^{PE}<p^I<p_0,\qquad n^D<n^I<n^F<1.
\]
\end{frame}

% ---------------------------------------------------------------------------
\begin{frame}{Alternative 3: slow cohort adjustment}
\begin{columns}[T,onlytextwidth]
\begin{column}{0.49\textwidth}
\centering
\textbf{Fertility}\par\smallskip
\includegraphics[width=\textwidth,height=0.51\textheight,keepaspectratio,trim=276 0 0 18,clip]{../output/model/fertility_population_housing_transition_note/steady_state_comparative_statics.pdf}
\end{column}
\begin{column}{0.49\textwidth}
\centering
\textbf{Housing market}\par\smallskip
\includegraphics[width=\textwidth,height=0.51\textheight,keepaspectratio,trim=0 0 276 18,clip]{../output/model/fertility_population_housing_transition_note/steady_state_comparative_statics.pdf}
\end{column}
\end{columns}

\vspace{-0.2em}
\small
\begin{itemize}
  \item At $I$, $n^I<1$: smaller cohorts gradually reduce housing demand and $p$.
  \item Lower $p$ raises fertility along the new curve until the economy reaches $A'$ with $n=1$ and $S_1<S_0$.
\end{itemize}
\end{frame}

% ---------------------------------------------------------------------------
\begin{frame}{The toy model now has a genuine transition}
After the permanent fall in the value of children, $v_0\rightarrow v_1<v_0$, each date is determined by
\[
\begin{aligned}
Y_{t+1}&=n(p_t;v_1)Y_t,
& O_{t+1}&=Y_t,\\
Y_t h_Y(p_t;v_1)+O_t h_O(p_t)&=H^s(p_t).
\end{aligned}
\]
The first line advances the two cohorts; the second clears the housing market given their current sizes.

\medskip
With total adult-household population $S_t=Y_t+O_t$,
\[
S_{t+1}-S_t
=Y_t\left(n_t-\frac{O_t}{Y_t}\right).
\]
Population therefore bottoms when fertility crosses the inherited old-to-young cohort ratio, not necessarily when fertility crosses one.

\smallskip
\centering
\large
$A\;\longrightarrow\;I\;\longrightarrow\;P\;\longrightarrow\;M\;\longrightarrow\;A'.$

\smallskip
\scriptsize
$A$: initial steady state; $I$: impact equilibrium; $P$: price trough; $M$: population trough; $A'$: new steady state.
\end{frame}

% ---------------------------------------------------------------------------
\begin{frame}{Population and housing costs undershoot the new steady state}
\centering
\includegraphics[width=0.97\textwidth,height=0.68\textheight,keepaspectratio]{../output/model/fertility_population_housing_transition_note/two_cohort_transition_paths.pdf}

\vspace{-0.2em}
\small
In this illustration, housing cost reaches its minimum at $P$ in generation 4 and population reaches its minimum at $M$ in generation 5. Both then recover slightly toward $A'$.
\end{frame}

% ---------------------------------------------------------------------------
\begin{frame}{The price trough comes just before the population trough}
\centering
\includegraphics[width=0.97\textwidth,height=0.66\textheight,keepaspectratio]{../output/model/fertility_population_housing_transition_note/two_cohort_trough_mechanics.pdf}

\vspace{-0.25em}
\small
Population keeps falling while $n_t<O_t/Y_t$, even after fertility has returned above one. The two troughs need not coincide because housing demand also depends on the composition of the population between young and old households.
\end{frame}

\end{document}
